\section{Discussion and Limitations}
\label{sec:discussion}

\paragraph{Validation scope.}
\sys validates enumeration, driver initialization, runtime control, DMA, and
interrupt handling against a physical backend. Because it terminates the
DUT-side accesses above the PCIe link, it cannot validate PHY training, LTSSM,
link errors, DLLP/TLP replay, equalization, or electrical behavior. It
complements rather than replaces a physical protocol analyzer or direct-link
test.

\paragraph{Extending device types.}
A new queue-based PCIe device can reuse topology, ECAM, dispatch, and DMA-window
mechanisms, but still needs a semantic adapter that identifies submission
points, translates its DMA-bearing structures, and defines completion ordering.
The NIC experience shows why generic MMIO forwarding is insufficient. A useful
future metric is adapter complexity and the fraction of logic reused across
devices, rather than device count alone.

\paragraph{Performance envelope.}
Software control mediation is appropriate when setup and submission rates are
modest relative to payload work. Faster DUTs, deep queues, or packet-scale
doorbells may expose the tens-of-microseconds path as a bottleneck. Batching,
doorbell coalescing, faster event transport, and selectively moving stable
parsing operations into parameterized hardware are possible responses without
abandoning the overall separation.

\paragraph{Topology and isolation.}
The current system configures a pre-provisioned hierarchy before boot and gives
each physical function to one backend. It does not provide hotplug, arbitrary
switch construction, SR-IOV sharing, malicious-DUT isolation, or live migration.
Those features require explicit IOMMU policy, reset containment, per-vector
interrupt remapping, and failure recovery.

\paragraph{Beyond pre-silicon prototypes.}
The logical subsystem abstraction could also front a fabricated DUT or support
user-space and paravirtual drivers. This produces an attractive longer-term
design space across development stage, driver model, and backend placement.
The present paper does not count those combinations as implemented features;
post-silicon reuse requires a compatible access front-end and separate evidence.

\paragraph{Interpretation of the contribution.}
The central claim is methodological: a stable, software-configurable device
presentation can be separated from backend realization when control semantics
and payload movement are treated as different planes and connected by explicit
coherence rules. The current prototype is evidence for that method in one
pre-silicon slice. MSI/MSI-X, multi-queue operation, AER, hotplug, Vortex/GDS,
post-silicon reuse, and the full $2\times3\times2$ space are extensions, not
unstated results of the present evaluation.
